\title{Aula 16 - Vulnerabilidade de Software}

\author{Prof. Gabriel Rodrigues Caldas de Aquino}

\institute
{
    Instituto de Computação \\
    Universidade Federal do Rio de Janeiro\\
    gabrielaquino@ic.ufrj.br% Your institution for the title page
}
\date{Compilado em: \\ \today} % Date, can be changed to a custom date

%----------------------------------------------------------------------------------------
%    PRESENTATION SLIDE
%----------------------------------------------------------------------------------------



\begin{frame}
    % Print the title page as the first slide
    \titlepage
\end{frame}


\begin{frame}{Segurança de Software e Programação Defensiva}
    \begin{itemize}
        \item \textbf{A Raiz do Problema:}
              \begin{itemize}
                  \item Muitas vulnerabilidades de segurança resultam de \textbf{práticas de programação ruins}.
                  \item A conscientização sobre essas falhas é um passo inicial crítico para escrever código de programa mais seguro.
                  \item A lista dos 10 Principais Riscos de Segurança em Aplicações Web da OWASP inclui cinco falhas diretamente relacionadas a código de software inseguro.

                  \item \textbf{Estas falhas incluem:}
                        \begin{itemize}
                            \item Entrada de dados não validada.
                            \item Cross-site scripting (XSS).
                            \item Estouro de buffer.
                            \item Falhas de injeção.
                            \item Tratamento inadequado de erros.
                        \end{itemize}
              \end{itemize}
    \end{itemize}
\end{frame}

\begin{frame}{As Falhas Mais Perigosas (CWE/SANS Top 25)}
    \begin{itemize}
        \item A lista \textbf{CWE/SANS Top 25 Most Dangerous Software Errors} detalha o consenso sobre práticas que causam a maioria dos ciberataques.

        \item \textbf{As falhas são agrupadas em 3 categorias:}
              \begin{itemize}
                  \item \textbf{1. Interação Insegura entre Componentes.}
                  \item \textbf{2. Gerenciamento de Recursos Arriscado.}
                  \item \textbf{3. Defesas Porosas.}
              \end{itemize}
    \end{itemize}
\end{frame}

\begin{frame}{Categoria de Erro de Software: Interação Insegura entre Componentes}
    \begin{itemize}
        \item Neutralização Imprópria de Elementos Especiais usados em um Comando SQL (``Injeção de SQL'')
        \item Neutralização Imprópria de Elementos Especiais usados em um Comando do SO (``Injeção de Comando do SO'')
        \item Neutralização Imprópria de Entrada Durante a Geração de Página Web (``Cross-site Scripting'')
        \item Upload Irrestrito de Arquivo com Tipo Perigoso
        \item Falsificação de Solicitação entre Sites (CSRF)
        \item Redirecionamento de URL para Site Não Confiável (``Redirecionamento Aberto'')
    \end{itemize}
\end{frame}

\begin{frame}{Categoria de Erro de Software: Gerenciamento de Recursos Arriscado}
    \begin{itemize}
        \item Cópia de Buffer sem Verificar o Tamanho da Entrada (``Estouro de Buffer Clássico'')
        \item Limitação Imprópria de um Nome de Caminho para um Diretório Restrito (``Path Traversal'')
        \item Download de Código Sem Verificação de Integridade
        \item Inclusão de Funcionalidade de Esfera de Controle Não Confiável
        \item Uso de Função Potencialmente Perigosa
        \item Cálculo Incorreto do Tamanho do Buffer
        \item String de Formatação Não Controlada
        \item Estouro de Inteiro ou \textit{Wraparound}
    \end{itemize}
\end{frame}